% id: 27-05
% book: 베이직쎈 미적분1 · 01 함수의 극한
% section: 학교 시험 기출로 실전 감각 UP!
% figure: none
% answer: 
% answer_source: 
% variant_level: 0 (원본 전사)
% 이 파일은 build.mjs 가 items.json 에서 만든다. 고칠 때는 items.json 을 고친다.
\smash{\raisebox{4.5mm}{\dmkichul{베이직쎈 미적분1 27쪽 05번}}}
함수의 극한에 대한 설명으로 옳은 것만을 \textbf{보기}에서 있는 대로 고른 것은? \bogi{ㄱ. $\displaystyle\lim_{x\to a} f(x)g(x)$의 값이 존재하면 $\displaystyle\lim_{x\to a} f(x)$, $\displaystyle\lim_{x\to a} g(x)$의 값도 각각 존재한다.\\ ㄴ. $\displaystyle\lim_{x\to a}\{f(x)-g(x)\}=0$이면 $\displaystyle\lim_{x\to a} f(x)=\lim_{x\to a} g(x)$이다.\\ ㄷ. $\displaystyle\lim_{x\to\infty} f(x)$와 $\displaystyle\lim_{x\to\infty}\dfrac{g(x)}{f(x)}$의 값이 모두 존재하면 $\displaystyle\lim_{x\to\infty} g(x)$의 값도 존재한다.}
\dmchoicesauto{}{ㄱ}{ㄷ}{ㄱ, ㄴ}{ㄴ, ㄷ}{ㄱ, ㄴ, ㄷ}
